\clearpage
\begin{center}
{\Large BESTIMMUNG}\par\vspace{8mm}
{\small DES}\par\vspace{8mm}
{\LARGE\textls[100]{BREITENUNTERSCHIEDES}}\par\vspace{8mm}
{\small ZWISCHEN DEN}\par\vspace{7mm}
{\large STERNWARTEN VON GÖTTINGEN UND ALTONA}\par\vspace{7mm}
{\small DURCH}\par\vspace{7mm}
{\normalsize BEOBACHTUNGEN AM RAMSDENSCHEN ZENITHSECTOR}\par\vspace{7mm}
{\small VON}\par\vspace{7mm}
{\large\textls[100]{CARL FRIEDRICH GAUSS,}}\par\smallskip
{\scriptsize
RITTER DES GUELPHEN- UND DANNEBROG-ORDENS; K. GROSSBR. HANNOVERSCHEM HOFRATH;\\
PROFESSOR DER ASTRONOMIE UND DIRECTOR DER STERNWARTE IN GÖTTINGEN;\\
MITGLIED DER AKADEMIEN UND SOCIETÄTEN VON BERLIN, COPENHAGEN, EDINBURG, GÖTTINGEN,\\
LONDON, MÜNCHEN, NEAPEL, PARIS, PETERSBURG, STOCKHOLM,\\
DER AMERIKANISCHEN, ITALIENISCHEN, KURLÄNDISCHEN, LONDONER ASTRONOMISCHEN U. A.
}\par\vspace{15mm}
\rule{.35\textwidth}{.4pt}\par\medskip
{\large\textls[170]{GÖTTINGEN,}}\par\medskip
BEI VANDENHOECK UND RUPRECHT.\par\medskip
1828.
\end{center}
\vfill\noindent IX.\hfill 1\par

\clearpage
\null

\clearpage
\begin{center}
{\LARGE\textls[170]{INHALT.}}\par\medskip
\rule{.09\textwidth}{.4pt}
\end{center}
\noindent Einleitung\dotfill S.\quad 5\par
\noindent\makebox[8mm][r]{I.}\quad Die beobachteten Sterne.\dotfill \quotedblbase\quad 8\par
\noindent\makebox[8mm][r]{II.}\quad Die Beobachtungen\dotfill \quotedblbase\quad 10\par
\noindent\makebox[8mm][r]{III.}\quad Resultate\dotfill \quotedblbase\quad 29\par
\noindent\hspace*{12mm}\begin{minipage}{\dimexpr\textwidth-24mm\relax}
Einfachste Combination der Beobachtungen zur Bestimmung des
Breitenunterschiedes, Art.\ 1. Genauigkeit der Beobachtungen,
Art.\ 2. Collimationsfehler, Art.\ 3. Absolut vortheilhafteste Combination
der Beobachtungen, Art.\ 4---7. Berücksichtigung der
unregelmässigen Theilungsfehler, Art.\ 8---11. Lage der Beobachtungsplätze,
Art.\ 12. Bestimmung der absoluten Polhöhe der
Göttinger Sternwarte aus Beobachtungen des Nordsterns am
\textsc{Reichenbach}schen Meridiankreise, Art.\ 13---17. Endresultat der
hannoverschen Gradmessung, Art.\ 18---20. Vergleichung der
Declinationen der beobachteten Zenithsterne mit \textsc{Bradleys} und
\textsc{Piazzis} Bestimmungen, Art.\ 21.
\end{minipage}\par
\noindent\makebox[8mm][r]{IV.}\quad Breitenbestimmung der Sternwarte Seeberg\dotfill \quotedblbase\quad 52\par
\noindent Zusatz zu [Art.\ 20] S.\ 48\dotfill \quotedblbase\quad 56\par
\begin{center}\rule{.23\textwidth}{.4pt}\end{center}
\vfill\begin{flushright}$1^{\ast}$\end{flushright}

\clearpage
\null

\clearpage
\vspace*{12mm}
\begin{center}
{\large BESTIMMUNG DES BREITENUNTERSCHIEDES}\par\medskip
{\small ZWISCHEN}\par\medskip
{\large DEN STERNWARTEN VON GÖTTINGEN UND ALTONA\\
DURCH BEOBACHTUNGEN\\
AM RAMSDENSCHEN ZENITHSECTOR.}\par\medskip
\rule{.12\textwidth}{.4pt}\par\bigskip
{\large EINLEITUNG.}
\end{center}

Durch die von mir in den Jahren 1821---1824 durch das Königreich Hannover
längs des Meridians von Göttingen geführte Dreieckskette sind die
Sternwarten von Göttingen und Altona auf das Genaueste trigonometrisch mit
einander verbunden. Diese Messungen werden in Zukunft ausführlich bekannt
gemacht werden: hier wird nur bemerkt, dass die absoluten Grössen
auf der von Herrn Prof.\ \textsc{Schumacher} in Holstein mit äusserster Schärfe gemessenen
Basis beruhen, mit welcher das Dreieckssystem durch die Seite
Hamburg-Hohenhorn zusammenhängt; die Orientirung gründet sich auf die
Beobachtungen am Göttinger Mittagsfernrohr, da die Sternwarte und das
nördliche Meridianzeichen selbst Dreieckspunkte sind. Die Sternwarten von
Göttingen und Altona liegen durch ein merkwürdiges Spiel des Zufalls auf
weniger als Eine Hausbreite in einerlei Meridian. Obgleich die absoluten
Polhöhen durch die Beobachtungen mit festen Meridianinstrumenten bestimmt
sind, so war es doch wichtig, den Unterschied der Breiten noch auf eine
andere Art mit einerlei Instrument zu bestimmen, und ich war so glücklich,

\clearpage
\noindent\makebox[10mm][l]{6}\parbox[t]{\dimexpr\textwidth-10mm\relax}{\centering\scriptsize BESTIMMUNG DES\,\quotedblbase\,BREITENUNTERSCHIEDES ZWISCHEN GÖTTINGEN UND ALTONA.}\par
\noindent dazu den trefflichen \textsc{Ramsden}schen Zenithsector anwenden zu können, der bekanntlich
zu ähnlichem Zweck bei der englischen Gradmessung gedient hat.
Die damit im Frühjahr 1827 von mir angestellten Beobachtungen und ihre
Resultate sind der Hauptgegenstand dieser Schrift.

Da die Beobachtungen mit diesem Instrument, wenn viele Sterne in einer
Reihe zu beobachten sind, nicht wohl ohne den Beistand eines geübten Gehülfen
gemacht werden können, so hatte Hr.\ Prof.\ \textsc{Schumacher} die Güte,
den Hrn.\ Ingenieur-Lieutenant v.\ \textsc{Nehus}[\textsuperscript{*})], unter Genehmigung Sr.\ Majestät
des Königs von Dänemark, für die Beobachtungen an beiden Plätzen damit zu
beauftragen. Dieser sehr geschickte Beobachter hat fortwährend die Ablesung
der Mikrometerschraube und die Einstellung des Lothfadens besorgt, während
ich selbst die Antritte an die Meridianfäden beobachtete und den auf den
Meridian senkrechten Faden auf die Sterne einstellte: nur in den beiden
ersten Beobachtungsnächten in Altona war jenes Geschäft von einem andern
Gehülfen besorgt; allein diese Beobachtungen sind deshalb nicht mit aufgenommen,
zumal da die Erfahrung bestätigte, dass verschiedene Personen die
Bisection der Punkte durch den Lothfaden ungleich schätzten.

Das Instrument ist durch die ausführliche von \textsc{Mudge} gegebene Beschreibung
hinlänglich bekannt. In Göttingen konnte es in der Sternwarte selbst,
unter dem östlichen Meridianspalt, aufgestellt werden. In Altona war dies
nicht thunlich; es wurde daher in dem Garten des Hrn.\ Prof.\ \textsc{Schumacher},
in welchem die dortige Sternwarte selbst liegt, unter demselben Beobachtungszelte,
welches \textsc{Mudge} in England gebraucht hat, aufgestellt. Die Solidität
der Aufstellung, auf eingerammten Pfählen, liess nichts zu wünschen übrig:
das Nivellement der Verticalaxe wurde täglich nachgesehen, und gewöhnlich
fast nichts zu ändern gefunden; dasselbe gilt von der Horizontalaxe.

Um die Ebene des Limbus in den Meridian zu bringen, wurde in Göttingen
das südliche Meridianzeichen benutzt, welches zwar in dem Meridian
des westlichen Spaltes steht, dessen Azimuth am Platze des Sectors sich aber
mit grösster Schärfe berechnen liess. In Altona konnte ein ähnliches Mittel
nicht angewandt werden: der Limbus wurde zuerst, mit Hülfe der Kenntniss
der absoluten Zeit, vermittelst eines culminirenden Sterns sehr nahe in den

\vfill\noindent\rule{.24\textwidth}{.4pt}\par
{\footnotesize\noindent [\textsuperscript{*}) Handschriftliche Bemerkung:] \textsc{Nehus} starb an zurückgetretener Grippe 1844 Apr.\ 17.\par}

\clearpage
\noindent\mbox{}\hfill EINLEITUNG.\hfill 7\par
\noindent Meridian gebracht; die Beobachtung mehrerer Sterne in der ganzen Ausdehnung
des Limbus gab dann leicht die noch nöthige kleine Correction der
Aufstellung. Da, wie schon erwähnt ist, von den am Sector culminirenden
Sternen in jeder Nacht auch die Antritte an die Meridianfäden beobachtet
wurden, und die Rectascensionen der Sterne bekannt waren, so erhielt die
richtige Aufstellung im Meridian dadurch eine fortwährende sichere Controlle,
und nur einmal war eine unbedeutende Nachhülfe erforderlich. In der Regel
wurde von einer Nacht zur andern mit der Stellung des Limbus, östlich oder
westlich, abgewechselt, und nur gegen den Schluss der Beobachtungen wurde,
um die Anzahl der Beobachtungen auf beide Lagen ziemlich gleich zu vertheilen,
von dieser Regel zuweilen abgewichen, und der Sector in Einer Nacht
ein oder mehreremale umgewandt.

Der Stand des Barometers und des innern und äussern Thermometers
wurde in jeder Nacht wenigstens dreimal, zu Anfang, in der Mitte und am
Schluss der Beobachtungen aufgezeichnet. Eben so, nach dem Vorgang von
\textsc{Mudge}, der Unterschied der Temperatur oben und unten am Sector, da deshalb
der Limbus und der Radius in ungleichem Verhältnisse verändert werden.
Dass übrigens jede andere durch die Einrichtung des Instruments vorgeschriebene
Vorsicht sorgfältig beachtet ist, z.\,B.\ das Wassergefäss, in welches das
Loth hängt, gehörig angefüllt zu erhalten, von der Mikrometerschraube so
viel thunlich dieselben Gewinde spielen zu lassen u.\ dergl., ist wohl überflüssig
besonders zu bemerken. Die Einstellung des Lothfadens auf den
obern Punkt (das Centrum des Gradbogens) wurde bei der Beobachtung jedes
Sterns von neuem unabhängig von der vorhergegangenen gemacht, und die
Einstellung auf den nächsten Theilungspunkt (oder die beiden nächsten) wurde
in der Regel mehreremale wiederholt, und aus den verschiedenen Ablesungen
der Mikrometerschraube, die meistens auf wenige Decimaltheile der Secunde
übereinstimmten, das Mittel genommen.

\begin{center}\rule{.24\textwidth}{.4pt}\end{center}

\clearpage
\vspace*{18mm}
\begin{center}
{\large I.}\par\medskip
{\Large DIE BEOBACHTETEN STERNE.}
\end{center}

Ich hatte zu Anfang 38 Sterne in schicklichen Lagen zur Beobachtung
ausgewählt, denen ich gegen den Schluss der Beobachtungen in Göttingen
noch fünf andere beifügte, weil ich besorgte, dass durch ungünstiges Wetter
der Schluss der Beobachtungen in Altona so weit verzögert werden könnte,
dass ein beträchtlicher Theil der ersten Sterne wegen der bei Tage eintretenden
Culmination nicht oft genug würde beobachtet werden können. Diese
Besorgniss bestätigte sich jedoch nur in geringem Grade, und nur ein einziger
Stern ist in Altona bloss einseitig beobachtet. Ich gebe hier die mittlere
Stellung dieser Sterne auf den Anfang des Jahrs 1827 reducirt: die Declinationen
sind die Resultate, welche die Beobachtungen am Zenithsector
selbst ergeben haben; die Rectascensionen, bei welchen für den gegenwärtigen
Zweck die allerschärfste Bestimmung unwesentlich ist, gründen sich meistens
nur auf eine einmalige Beobachtung am Meridiankreise, deren Reduction
Hr.\ \textsc{von Heiligenstein} gefälligst berechnet hat. Der Bequemlichkeit wegen
bezeichne ich die Sterne mit fortlaufenden Zahlen. Nro.\ 8, 13, 15 und 31
sind Doppelsterne; bei dem ersten ist immer der nachfolgende Stern, bei den
beiden folgenden die Mitte eingestellt; bei Nro.\ 31 ist der Nebenstern so
klein, dass er im Fernrohr des Sectors, selbst bei versuchsweise verdunkeltem
Felde, immer unsichtbar blieb, obgleich er von Hrn.\ Prof.\ \textsc{Schumacher} im
lichtstärkern Fernrohr des \textsc{Reichenbach}schen Meridiankreises sofort bemerkt
wurde, ohne dass uns damals bekannt war, dass schon andere Astronomen
diesen Doppelstern als solchen erkannt hatten.

\clearpage
\noindent\parbox[t]{\dimexpr\textwidth-10mm\relax}{\centering\scriptsize BESTIMMUNG DES BREITENUNTERSCHIEDES ETC.\quad I.\quad DIE BEOBACHTETEN STERNE.}\makebox[10mm][r]{9}\par
\begin{center}
\begin{minipage}[t]{.49\textwidth}\centering
{\scriptsize\setlength{\tabcolsep}{1.2pt}\renewcommand{\arraystretch}{1.2}\begin{tabular}{r|l|rrr|rrr|}
\multicolumn{2}{c|}{Bezeichnung}&\multicolumn{3}{c|}{G. Aufst. 1827}&\multicolumn{3}{c|}{Declin. 1827}\\\hline
1&24 Canum&$13^{\mathrm h}$&$27^{\mathrm m}$&$22^{\mathrm s}{,}39$&$49^\circ$&$54'$&$11''{,}62$\\
2&83 Ursae&&$34$&$9{,}77$&$55$&$33$&$34{,}98$\\
3&$\eta$ Ursae&&$40$&$42{,}86$&$50$&$10$&$46{,}20$\\
4&86 Ursae&&$47$&$28{,}48$&$54$&$34$&$55{,}55$\\
5&&&$50$&$45{,}55$&$55$&$25$&$59{,}22$\\
6&P. 13. 289&&$55$&$19{,}44$&$46$&$35$&$39{,}14$\\
7&13 Bootis&$14$&$1$&$49{,}02$&$50$&$16$&$43{,}44$\\
8&$\varkappa$ Bootis seq.&&$7$&$16{,}93$&$52$&$36$&$7{,}47$\\
9&P. 14. 56&&$12$&$5{,}39$&$56$&$13$&$36{,}51$\\
10&$\theta$ Bootis&&$19$&$18{,}40$&$52$&$39$&$12{,}05$\\
11&P. 14. 131&&$27$&$50{,}42$&$53$&$39$&$33{,}02$\\
12&P. 14. 164&&$35$&$24{,}18$&$52$&$58$&$55{,}66$\\
13&39 Bootis med.&&$43$&$48{,}37$&$49$&$26$&$9{,}46$\\
14&P. 14. 235&&$50$&$38{,}71$&$50$&$20$&$22{,}49$\\
15&44 Bootis med.&&$58$&$5{,}30$&$48$&$19$&$52{,}47$\\
16&&$15$&$7$&$8{,}12$&$49$&$13$&$46{,}68$\\
17&P. 15. 39&&$10$&$33{,}95$&$51$&$34$&$53{,}72$\\
18&&&$15$&$0{,}55$&$52$&$35$&$5{,}87$\\
19&&&$21$&$50{,}88$&$54$&$37$&$34{,}87$\\
20&&&$30$&$42{,}13$&$54$&$29$&$54{,}30$\\
21&&&$38$&$9{,}07$&$52$&$54$&$37{,}09$\\
22&&&$42$&$14{,}32$&$46$&$16$&$4{,}70$\\
\end{tabular}}\end{minipage}\hfill
\begin{minipage}[t]{.49\textwidth}\centering
{\scriptsize\setlength{\tabcolsep}{1.2pt}\renewcommand{\arraystretch}{1.2}\begin{tabular}{r|l|rrr|rrr}
\multicolumn{2}{c|}{Bezeichnung}&\multicolumn{3}{c|}{G. Aufst. 1827}&\multicolumn{3}{c}{Declin. 1827}\\\hline
23&&$15^{\mathrm h}$&$48^{\mathrm m}$&$16^{\mathrm s}{,}24$&$56^\circ$&$20'$&$27''{,}00$\\
24&&&$54$&$10{,}23$&$50$&$22$&$39{,}41$\\
25&$\theta$ Draconis&&$58$&$39{,}70$&$59$&$1$&$47{,}14$\\
26&&$16$&$4$&$24{,}06$&$50$&$38$&$13{,}12$\\
27&P. 16. 33&&$7$&$24{,}92$&$46$&$20$&$17{,}04$\\
28&P. 16. 56&&$11$&$34{,}31$&$53$&$40$&$13{,}41$\\
29&&&$16$&$2{,}19$&$52$&$27$&$11{,}06$\\
30&&&$20$&$38{,}75$&$55$&$36$&$4{,}73$\\
31&&&$26$&$34{,}90$&$45$&$58$&$5{,}85$\\
32&16 Draconis&&$32$&$6{,}42$&$53$&$15$&$3{,}64$\\
33&&&$37$&$53{,}85$&$50$&$16$&$13{,}67$\\
34&&&$42$&$1{,}51$&$57$&$5$&$37{,}79$\\
35&&&$45$&$2{,}26$&$46$&$56$&$48{,}53$\\
36&P. 16. 253&&$49$&$21{,}51$&$46$&$49$&$22{,}01$\\
37&P. 16. 291&&$56$&$11{,}73$&$56$&$56$&$42{,}35$\\
38&P. 16. 310&$17$&$0$&$15{,}10$&$49$&$2$&$47{,}09$\\
39&P. 17. 20&&$4$&$23{,}12$&$58$&$29$&$49{,}33$\\
40&P. 17. 38&&$8$&$1{,}46$&$56$&$52$&$21{,}74$\\
41&74 Herculis&&$15$&$28{,}30$&$46$&$24$&$52{,}90$\\
42&P. 17. 120&&$20$&$53{,}56$&$57$&$10$&$13{,}11$\\
43&$\beta$ Draconis&&$26$&$31{,}90$&$52$&$25$&$57{,}91$\\
&&&&&&&\\
\end{tabular}}\end{minipage}
\end{center}
\begin{center}\rule{.24\textwidth}{.4pt}\end{center}

\clearpage
\vspace*{18mm}
\begin{center}
{\large II.}\par\medskip
{\Large DIE BEOBACHTUNGEN.}
\end{center}

Ein vollständiger Abdruck des Tagebuchs in seiner ursprünglichen Gestalt,
welcher die Stärke dieser Schrift mehr als verdoppelt haben würde, hat
mir überflüssig geschienen: ich gebe daher die Beobachtungen sogleich nach
den Sternen geordnet.

Die erste Columne enthält die Zenithdistanzen, wie das Instrument sie
gegeben hat, d.\,i.\ die blosse Reduction der Ablesung. Nördliche Zenithdistanzen
sind als positiv, südliche als negativ betrachtet.

Die zweite Columne gibt die Vereinigung der Refraction mit der Wirkung
der ungleichen Ausdehnung des Instruments wegen Ungleichheit der
obern und untern Temperatur: die äussersten vorgekommenen Unterschiede
waren $+1^\circ{,}2$ Réaum.\ (das obere Thermometer höher) und $-0^\circ{,}6$. Um ein
reines Resultat zur Beurtheilung der Übereinstimmung der Beobachtungen
unter sich zu erhalten, habe ich die Mühe nicht gescheut, den Betrag für
jede einzelne Beobachtung zu berechnen, wobei jedoch einige kleine sich
leicht darbietende Rechnungsvortheile benutzt sind.

Die dritte Columne enthält die Reduction auf den mittlern Ort für den
Anfang des Jahrs wegen Aberration, Nutation und Praecession, wozu bei
einigen Sternen noch die eigene Bewegung gesetzt ist: es ist nemlich die jährliche
eigene Bewegung in Declination angenommen
\begin{center}
\begin{tabular}{@{}lr@{\quad}c@{\quad}l@{}}
für&10&.\ .\ .\ .\ .\ .\ .&$-0''{,}42$\\
&25&.\ .\ .\ .\ .\ .\ .&$+0{,}33$\\
&37&.\ .\ .\ .\ .\ .\ .&$+0{,}38$.
\end{tabular}
\end{center}

\clearpage
\noindent\parbox[t]{\dimexpr\textwidth-10mm\relax}{\centering\scriptsize BESTIMMUNG DES BREITENUNTERSCHIEDES ETC.\quad II.\quad DIE BEOBACHTUNGEN.}\makebox[10mm][r]{11}\par
\noindent Bei den beiden ersten Sternen ist die eigene Bewegung längst entschieden;
bei 37 zeigt sie sich durch Vergleichung mit \textsc{Piazzis} Bestimmung so, dass,
die Richtigkeit der letztern vorausgesetzt, sie nicht bezweifelt werden kann\textsuperscript{*}).
Der Berechnung der Aberration, Nutation und Praecession liegen \textsc{Bailys}
schätzbare Tafeln zum Grunde, nach denen für jeden Stern eine Ephemeride
von 10 zu 10 Tagen, unter Beihülfe der Hrn.\ v.\ \textsc{Nehus} und \textsc{Petersen}, berechnet,
und in diese mit Berücksichtigung der zweiten Differenzen interpolirt
wurde.

Die vierte Columne enthält endlich die Summe der drei ersten, also die
wahre nur noch mit dem Collimationsfehler behaftete Zenithdistanz für die
mittlere Stellung zu Anfang des Jahrs 1827, wie sie sich aus jeder einzelnen
Beobachtung ergibt.

\medskip\noindent\rule{.24\textwidth}{.4pt}\par
{\footnotesize
\noindent\textsuperscript{*}) Die Richtigkeit von \textsc{Piazzis} Bestimmung dieses Sterns, auf 8 Beobachtungen gegründet, erhält
durch die nahe Übereinstimmung mit der Angabe der ältern Ausgabe seines Verzeichnisses von 1803,
welche auf 6 Beobachtungen beruhte, eine Bestätigung; die genaue Grösse der eigenen Bewegung bleibt
aber deswegen noch etwas ungewiss, weil das Jahr unbekannt ist, welches dem Mittel der Beobachtungen
entspricht. Es ist merkwürdig, diese nicht unbeträchtliche eigene Bewegung bei einem Stern der 7.\ Grösse
zu finden. Auch Nro.\ 11 scheint in dieser Beziehung die Aufmerksamkeit der Astronomen zu verdienen.\par}
\begin{center}\rule{.24\textwidth}{.4pt}\end{center}
